% VI21-FULL-SOLUTION-REFINEMENT: P02
\subsection*{VI.21.P2: The affine formula for the structure morphism}
\begin{problem}\label{prob:vi21-p02}
Let $X=\Spec A$.  Show that $\pi_X$ is induced by $\mathbb Z\to A$ and prove
\[
\pi_X(\mathfrak p)=\mathfrak p\cap\mathbb Z
\]
for every $\mathfrak p\in\Spec A$.
\end{problem}

\ifdefined\FullProblemDossiers
\paragraph{Why this problem matters.}
It translates the final-object theorem into the concrete contraction formula used in arithmetic examples.

\paragraph{Useful ideas before starting.}
Translate point statements through local rings and residue fields; when the target is affine, use the VI/19 universal property and then verify the pointwise interpretation.

\paragraph{Guided hints.}
\begin{enumerate}
\item Identify the relevant canonical ring or stalk homomorphism.
\item Track maximal ideals or prime contractions to identify the image point.
\item Use locality to pass to residue fields when needed.
\item Check independence, uniqueness, or functoriality explicitly rather than only on points.
\end{enumerate}

\begin{solution}
By affine anti-equivalence, a morphism $\Spec A\to\Spec\mathbb Z$ corresponds to a unital ring map
$\mathbb Z\to A$.  There is only one such map, $n\mapsto n1_A$, so it must correspond to the
absolute structure morphism $\pi_X$.  The underlying point map of $\Spec(\mathbb Z\to A)$ is
contraction of prime ideals.  Hence
\[
\pi_X(\mathfrak p)=(\mathbb Z\to A)^{-1}(\mathfrak p)=\mathfrak p\cap\mathbb Z.
\]
Because the only prime ideals of $\mathbb Z$ are $(0)$ and $(p)$ for prime integers $p$, the
formula immediately records the characteristic of $\kappa(\mathfrak p)$.
\end{solution}

\paragraph{What happened mathematically.}
It translates the final-object theorem into the concrete contraction formula used in arithmetic examples.

\paragraph{Extensions.}
Apply it to $\Spec\mathbb Z[t]$, $\Spec\mathbb F_p[t]$, and mixed-characteristic rings.

\paragraph{Provenance.}
\textbf{Canonical VI/21 derivative/application; not an additional legacy migration.}
\else
\paragraph{Provenance.} \textbf{Canonical VI/21 derivative/application; not an additional legacy migration.}
\fi
